% =====================================================================
\chapter{Lua-скрипты}
\label{ch:lua}
% =====================================================================
\chapterintro
  {какие бывают Lua-скрипты в EdgeTX, где они лежат и как запускаются; готовые полезные
   скрипты (ExpressLRS, Betaflight); как написать свой экран телеметрии, микшерный и
   функциональный скрипт}
  {SD-карта пульта, текстовый редактор на компьютере}

\section{Зачем Lua}

\textbf{Lua} --- небольшой встраиваемый язык программирования. EdgeTX умеет выполнять
Lua-скрипты с SD-карты: так работают настройщик ExpressLRS, конфигуратор Betaflight прямо
с пульта, мастера настройки моделей, экраны телеметрии, игры. Если какой-то логики не
хватает в стандартных меню --- её можно дописать на Lua.

\section{Типы скриптов}

\begin{center}\small
\renewcommand{\arraystretch}{1.2}
\begin{tabularx}{\linewidth}{@{}p{28mm}p{38mm}X@{}}
\toprule
Тип & Папка & Как запускается \\
\midrule
Утилиты (Tools, one-time) & \code{SCRIPTS/TOOLS} & \mpath{SYS\sep Tools}; занимают экран,
  пока не завершатся \\
Экраны телеметрии & \code{SCRIPTS/TELEMETRY} & \mpath{MDL\sep Telemetry\sep Screens\sep Script};
  показываются по \btn{TELE} \\
Функциональные & \code{SCRIPTS/FUNCTIONS} & специальная функция \menu{Lua script}
  работает, пока условие активно \\
Микшерные & \code{SCRIPTS/MIXES} & \mpath{MDL\sep Custom Scripts}; выходы скрипта становятся
  источниками для миксеров \\
Мастера (wizards) & \code{SCRIPTS/WIZARD} & при создании модели \\
\bottomrule
\end{tabularx}
\end{center}

Виджеты (\code{WIDGETS}) работают только на пультах с цветным экраном и на TX12/Boxer
не используются.

\section{Полезные готовые скрипты}

\begin{description}[font=\sffamily\bfseries\color{etx}]
  \item[ExpressLRS (\code{elrsV3.lua})] настройка модуля и приёмника ELRS
        (глава~\ref{ch:rf}). Через пункт \menu{Other Devices} можно настроить и сам
        приёмник, а через \menu{VTX Administrator} --- канал видеопередатчика.
  \item[Betaflight TX Lua Scripts] настройка PID, рейтов, фильтров, VTX на квадрокоптере с
        Betaflight прямо с пульта через телеметрию CRSF/SmartPort. Скачиваются с GitHub
        проекта Betaflight, копируются в \code{SCRIPTS}, запускаются из \menu{Tools}.
  \item[Мастера моделей] входят в пакет SD-карты EdgeTX.
\end{description}

\section{Устройство скрипта}

Любой скрипт EdgeTX --- это файл, который \emph{возвращает таблицу} с функциями:
\begin{itemize}
  \item \code{init()} --- вызывается один раз при загрузке;
  \item \code{run(...)} --- вызывается многократно: для утилит и экранов --- с кодом
        нажатой кнопки (\code{event}), для микшерных --- со значениями входов;
  \item \code{background()} --- (для экранов телеметрии) вызывается, пока экран не показан.
\end{itemize}

Ключевые функции API:
\begin{center}\small
\begin{tabularx}{\linewidth}{@{}p{55mm}X@{}}
\toprule
\code{getValue("RxBt")} & значение источника/датчика по имени (0, если нет) \\
\code{lcd.clear()} & очистить экран \\
\code{lcd.drawText(x, y, s, flags)} & текст; флаги \code{INVERS}, \code{BLINK},
  \code{SMLSIZE}, \code{MIDSIZE}, \code{DBLSIZE} \\
\code{lcd.drawNumber(x, y, n, flags)} & число; \code{PREC1}/\code{PREC2} --- 1 или 2 знака
  после запятой (передавайте $n \times 10$ или $n \times 100$) \\
\code{lcd.drawRectangle}, \code{lcd.drawFilledRectangle}, \code{lcd.drawLine} & графика \\
\code{model.getInfo()}, \code{model.getTimer(i)} & данные текущей модели \\
\code{playNumber(n, unit, flags)}, \code{playFile(path)} & голос \\
\code{getTime()} & время в «тиках» по 10\,мс \\
\bottomrule
\end{tabularx}
\end{center}

Экран монохромный, 128×64 точки: $x$ от 0 до 127 слева направо, $y$ от 0 до 63 сверху вниз.

\section{Пример 1: утилита}

\luafile{hello.lua}

Скопируйте файл в \code{SCRIPTS/TOOLS/hello.lua}, откройте \mpath{SYS\sep Tools} и
запустите \menu{hello}. \code{return 1} из \code{run} завершает утилиту.

\section{Пример 2: экран телеметрии}

\luafile{batt.lua}

Файл кладётся в \code{SCRIPTS/TELEMETRY/batt.lua}. На монохромных пультах имя файла
телеметрического скрипта должно быть коротким --- не длиннее 6 символов без расширения.
Затем \mpath{MDL\sep Telemetry\sep Screen 1} = \menu{Script}, \menu{batt}.

\section{Пример 3: микшерный скрипт}

\luafile{thrlim.lua}

Скрипт получает стик газа и параметр \code{Limit} и отдаёт ограниченный газ. В
\mpath{MDL\sep Custom Scripts} выберите \menu{thrlim}, укажите источник \sw{Thr} и предел.
В миксере канала газа используйте источник \sw{Lim}. Микшерные скрипты выполняются
каждый цикл, поэтому они должны быть быстрыми; для управления опасными функциями
(моторы) применяйте их с осторожностью: при ошибке в скрипте EdgeTX его отключит, и выход
перестанет обновляться.

\section{Пример 4: функциональный скрипт}

\luafile{lqcall.lua}

Кладётся в \code{SCRIPTS/FUNCTIONS/lqcall.lua}. Специальная функция: \menu{Switch} = \sw{SA↓},
\menu{Lua script} = \menu{lqcall}. В арме пульт раз в 10 секунд называет качество связи.

\section{Отладка скриптов}

\begin{itemize}
  \item Синтаксические ошибки EdgeTX показывает на экране при запуске (\code{Script syntax error}).
  \item Проверяйте скрипты на компьютере: EdgeTX Companion содержит симулятор пульта с
        поддержкой Lua и отладочным выводом \code{print()}.
  \item Если скрипт \emph{экрана} или \emph{миксера} «съел» слишком много времени или памяти,
        EdgeTX его остановит (\code{Script killed}). Упрощайте код, не создавайте таблицы в
        каждом вызове \code{run}.
\end{itemize}

\begin{tasks}
\begin{enumerate}
  \item Запустите \code{hello.lua} и добавьте в него вывод значения таймера 1.
  \item Настройте \code{batt.lua} под свой аккумулятор (\code{CELLS}) и выведите его по \btn{TELE}.
  \item Измените \code{lqcall.lua}, чтобы он говорил напряжение \sw{RxBt}, а не \sw{RQly}.
\end{enumerate}
\end{tasks}
